\setcounter{section}{11}

  \zwischenueberschrift{Produk manifold}

\inputdefinition
{}
{

Misalkan
\mathkor {} {M} {dan} {N} {}
dua
\definitionsverweis {manifold diferensiabel}{}{}
dengan
\definitionsverweis {atlas-atlas}{}{}
\mathkor {} {(U_i,U_i',\alpha_i, i \in I)} {dan} {(V_j,V_j',\beta_j, j \in J)} {.}
Maka
\definitionsverweis {ruang produk}{}{}

\mathl{M \times N}{} dengan
\definitionsverweis {peta-peta}{}{}
\maabbdisp {\alpha_i \times \beta_j} {U_i \times V_j } { U_i' \times V_j'
} {}
\zusatzklammer {dengan

\mavergleichskettek
{\vergleichskettek
{  (i,j)
}
{ \in }{ I \times J
}
{  }{
}
{    }{
}
{   }{
}
}
{}{}{}
dan

\mavergleichskettek
{\vergleichskettek
{ U_i' \times V_j'
}
{ \subseteq }{ \R^m \times \R^n
}
{  }{
}
{    }{
}
{   }{
}
}
{}{}{}} {} {}
disebut \definitionswort {produk manifold}{}
\mathkor {} {M} {dan} {N} {.}

}

Objek ini memang merupakan suatu manifold; lihat
Soal 11.1.
Manifold produk berbentuk
\mathl{M \times \R}{} disebut juga \stichwort {silinder} {} di atas $M$. Jika

\mavergleichskette
{\vergleichskette
{ M
}
{ \subseteq }{ \R^\ell
}
{  }{
}
{    }{
}
{   }{
}
}
{}{}{}
dan

\mavergleichskette
{\vergleichskette
{ N
}
{ \subseteq }{ \R^k
}
{  }{
}
{    }{
}
{   }{
}
}
{}{}{}
merupakan submanifold-submanifold tertutup, maka manifold produk
\mathl{M \times N}{} merupakan submanifold tertutup dari

\mavergleichskette
{\vergleichskette
{ \R^\ell \times \R^k
}
{ = }{ \R^{\ell + k}
}
{  }{
}
{    }{
}
{   }{
}
}
{}{}{,}
lihat
Soal 11.2.



\inputfaktbeweis
{Produkt von Mannigfaltigkeiten/Abbildungseigenschaften/Fakt}
{Lema}
{}
{

\faktsituation {Misalkan
\mathkor {} {M} {dan} {N} {}
\definitionsverweis {manifold-manifold diferensiabel}{}{}
dan
\mathl{M \times N}{} merupakan
\definitionsverweis {produk}{}{} keduanya.}
\faktuebergang {Maka berlaku sifat-sifat berikut.}
\faktfolgerung {\aufzaehlungdrei{Kedua
\definitionsverweis {proyeksi}{}{}
\maabbeledisp {p_1} {M \times N} {M
} {(x,y)} { x
} {,}
dan
\maabbeledisp {p_2} {M \times N} {N
} {(x,y)} { y
} {,}
merupakan
\definitionsverweis {pemetaan diferensiabel}{}{.}
}{
\definitionsverweis {Ruang singgung}{}{}
di suatu titik

\mavergleichskette
{\vergleichskette
{ R
}
{ = }{ (P,Q)
}
{  }{
}
{    }{
}
{   }{
}
}
{}{}{}
secara kanonik teridentifikasi dengan

\mavergleichskette
{\vergleichskette
{ T_R(M \times N)
}
{ = }{ T_PM \times T_QN
}
{  }{
}
{    }{
}
{   }{
}
}
{}{}{.}
}{Misalkan $L$ suatu manifold diferensiabel lain. Maka suatu pemetaan
\maabbeledisp {\varphi \times \psi} { L } { M \times N
} { u } { (\varphi(u), \psi(u))
} {,}
diferensiabel jika dan hanya jika kedua
\definitionsverweis {pemetaan komponennya}{}{}
\mathkor {} {\varphi} {dan} {\psi} {}
diferensiabel.
}}
\faktzusatz {}
\faktzusatz {}

}
{

\teilbeweis {}{}{}
{(1). Dengan beralih ke peta, kita dapat mengasumsikan bahwa
\mathkor {} {M} {dan} {N} {}
merupakan
\definitionsverweis {himpunan-himpunan bagian terbuka}{}{}
masing-masing di
\mathkor {} {\R^m} {dan} {\R^n} {.}
Dalam hal ini, pemetaan tersebut merupakan suatu
\definitionsverweis {pembatasan}{}{}
dari
\definitionsverweis {proyeksi linear}{}{}
\maabb {} {\R^m \times \R^n } {\R^m
} {,}
yang
menurut Proposisi 45.3 (Analysis (Osnabrück 2021-2023))
\definitionsverweis {terdiferensialkan secara kontinu}{}{.}}
{}
\teilbeweis {}{}{}
{(2). Proyeksi-proyeksi diferensiabel
\maabb {p_1} {M \times N} {M
} {}
dan
\maabb {p_2} {M \times N} {N
} {}
menghasilkan
\definitionsverweis {pemetaan-pemetaan singgung}{}{}
linear
\maabb {T_R(p_1)} {T_R(M \times N)} {T_PM
} {}
dan
\maabb {T_R(p_2)} {T_R(M \times N)} {T_QN
} {}
sehingga secara keseluruhan menghasilkan pemetaan linear
\maabbdisp {T_R(p_1) \times T_R(p_2)} { T_R(M \times N)} {T_PM \times T_QN} {.}
Untuk membuktikan bijektivitas, kita dapat beralih ke peta dan mengasumsikan bahwa

\mavergleichskette
{\vergleichskette
{ M
}
{ \subseteq }{ \R^m
}
{  }{
}
{    }{
}
{   }{
}
}
{}{}{}
dan

\mavergleichskette
{\vergleichskette
{ N
}
{ \subseteq }{ \R^n
}
{  }{
}
{    }{
}
{   }{
}
}
{}{}{}
merupakan himpunan-himpunan bagian terbuka. Pemetaan ini kemudian menjadi bijeksi
\maabbdisp {p_1 \times p_2} {\R^{m+n} } { \R^m \times \R^n
} {.}}
{}
\teilbeweis {}{}{}
{(3). Untuk suatu titik tetap

\mavergleichskette
{\vergleichskette
{ A
}
{ \in }{ L
}
{  }{
}
{    }{
}
{   }{
}
}
{}{}{}
dengan menggunakan lingkungan peta dari $A$ dan dari
\mathkor {} {\varphi(A)} {dan} {\psi(A)} {}
kita dapat mereduksi keadaan ke kasus ketika ketiga manifold merupakan himpunan-himpunan terbuka dalam ruang-ruang Euklides. Jika kedua pemetaan terdiferensialkan secara kontinu, maka menurut
Soal 45.9 (Analysis (Osnabrück 2021-2023))
pemetaan keseluruhannya terdiferensialkan secara kontinu. Arah sebaliknya jelas.}
{}

}

\bild{ \begin{center}
\includegraphics[width=5.5cm]{ \bildeinlesungpng {Toroidal_coord} {png} }
\end{center}
\bildtext {} }

\bildlizenz { Toroidal coord.png } {} {DaveBurke} {Commons} {CC-by 2.5} {}

\inputbeispiel{}
{

\definitionsverweis {Produk}{}{}
dari
\definitionsverweis {lingkaran}{}{}
dengan dirinya sendiri, yaitu

\mavergleichskette
{\vergleichskette
{ M
}
{ = }{ S^1 \times S^1
}
{  }{
}
{    }{
}
{   }{
}
}
{}{}{,}
disebut \stichwort {torus} {.} Ini merupakan manifold berdimensi dua. Karena

\mavergleichskette
{\vergleichskette
{ S^1
}
{ \subset }{ \R^2
}
{  }{
}
{    }{
}
{   }{
}
}
{}{}{}
merupakan
\definitionsverweis {submanifold tertutup}{}{,}
torus dapat direalisasikan sebagai submanifold tertutup di

\mavergleichskette
{\vergleichskette
{ \R^2 \times \R^2
}
{ = }{ \R^4
}
{  }{
}
{    }{
}
{   }{
}
}
{}{}{.}
Namun, torus juga dapat direalisasikan sebagai submanifold tertutup di $\R^3$. Untuk itu, misalkan
\mathkor {} {r} {dan} {R} {}
bilangan-bilangan real positif dengan

\mavergleichskette
{\vergleichskette
{ 0
}
{ < }{ r
}
{ < }{ R
}
{    }{
}
{   }{
}
}
{}{}{.}
Maka himpunan

\mathdisp {\left\{ (x,y,z) \in \R^3  \mid   \left(  \sqrt{x^2+y^2} -R  \right)^2 +z^2 = r^2         \right\}} {  }

merupakan suatu torus. Dalam realisasi ini, torus merupakan permukaan
\zusatzklammer {yang digelembungkan} {} {}
\anfuehrung{ban dalam sepeda}{,} dengan \anfuehrung{jari-jari roda}{} sama dengan $R$ dan \anfuehrung{jari-jari tabung}{} sama dengan $r$
\zusatzklammer {roda terletak pada bidang \mathlk{x-y}{-}} {} {.}
Setelah setiap faktor lingkaran diidentifikasi melalui
\mathl{S^1 \cong \R/(2\pi\mathbb{Z})}{,}
hubungan dengan produk
\mathl{S^1 \times S^1}{} diperoleh dengan memetakan pasangan kelas sudut
\mathl{([ \varphi ],[ \psi ])}{} ke titik
\mathl{( (R +r \cos  \psi) \cos \varphi, (R+ r \cos  \psi ) \sin \varphi ,  r \sin  \psi )}{}.

}

  \zwischenueberschrift{Braket Lie medan vektor}

\inputdefinition
{}
{

Misalkan $M$ suatu
\definitionsverweis {manifold diferensiabel}{}{.}
Suatu pemetaan
\maabbdisp {D} {C^1(M,\R) } { C^0(M,\R)
} {}
disebut
\definitionswort {operator diferensial orde pertama}{,}
jika memenuhi sifat-sifat berikut.
\aufzaehlungzwei { $D$ bersifat
$\R$-\definitionsverweis {linear}{}{.}
} {Berlaku

\mavergleichskette
{\vergleichskette
{ D(fg)
}
{ = }{ fD(g)+gD(f)
}
{  }{
}
{    }{
}
{   }{
}
}
{}{}{.}
}

}

Operator diferensial orde pertama disebut juga \stichwort {derivasi} {.}

Untuk suatu himpunan terbuka

\mavergleichskette
{\vergleichskette
{ U
}
{ \subseteq }{ \R^n
}
{  }{
}
{    }{
}
{   }{
}
}
{}{}{,}
turunan-turunan parsial $\partial_i$ dan $\partial_j$, sebagai operasi pada fungsi-fungsi yang dua kali terdiferensialkan secara kontinu, dapat dipertukarkan; jadi

\mavergleichskettedisp
{\vergleichskette
{ \partial_i \circ \partial_j
}
{ =} { \partial_j \circ \partial_i
}
{ } {
}
{ } {
}
{ } {
}
}
{}{}{.}
Suatu kombinasi linear

\mavergleichskettedisp
{\vergleichskette
{ V
}
{ =} { \sum_{i = 1}^n f_i \partial_i
}
{ } {
}
{ } {
}
{ } {
}
}
{}{}{}
dengan koefisien-koefisien kontinu
\mathl{f_i \in C^0(U,\R)}{}
dapat dipandang sebagai
\definitionsverweis {medan vektor}{}{}
\maabbeledisp {} { U } { TU = U \times \R^n
} { P } { (P, \sum_{i = 1}^n f_i(P) \partial_i ) = (P,f_1(P) , \ldots , f_n(P) )
} {,}
dan juga sebagai operator turunan
\maabbeledisp {} {C^1(U,\R)} {C^0(U,\R)
} {h } { \sum_{i = 1}^n f_i  \partial_i  h
} {,}
\zusatzklammer {atau dari $C^\infty(U,\R)$ ke $C^\infty(U,\R)$ apabila
\mathl{f_i \in C^\infty(U,\R)}{} untuk setiap $i$} {} {.}
Dalam pengertian ini, medan vektor dan operator diferensial orde pertama saling bersesuaian. Lebih lanjut, operator-operator diferensial dapat dikomposisikan satu sama lain, dan ternyata urutannya penting.



\inputfaktbeweis
{Offene Menge/R^n/Vektorfelder/Hintereinanderschaltung/Explizit/Fakt}
{Lema}
{}
{

\faktsituation {Misalkan

\mavergleichskette
{\vergleichskette
{ U
}
{ \subseteq }{ \R^n
}
{  }{
}
{    }{
}
{   }{
}
}
{}{}{}
\definitionsverweis {terbuka}{}{}
dan misalkan $f_1 , \ldots , f_n, g_1 , \ldots , g_n,h$ merupakan fungsi-fungsi pada $U$ yang dua kali
\definitionsverweis {terdiferensialkan secara kontinu}{}{.}}
\faktfolgerung {Maka berlaku

\mavergleichskettedisphandlinks
{\vergleichskettedisphandlinks
{ \left(  \sum_{i = 1}^n f_i \partial_i  \right) \left(  \left(  \sum_{j = 1}^n g_j \partial_j   \right) (h)  \right)
}
{ =} { \sum_{j = 1}^n  \left(  \sum_{i = 1}^n  f_i \partial_i g_j  \right) \partial_j h + \sum_{j = 1}^n \sum_{i = 1}^n    f_i g_j \partial_i \partial_j h
}
{ } {
}
{ } {
}
{ } {
}
}
{}{}{.}}
\faktzusatz {}
\faktzusatz {}

}
{

Berlaku

\mavergleichskettealignhandlinks
{\vergleichskettealignhandlinks
{ \left(  \sum_{i = 1}^n f_i \partial_i  \right) \left(  \left(  \sum_{j = 1}^n g_j \partial_j   \right) (h)  \right)
}
{ =} { \left(  \sum_{i = 1}^n f_i \partial_i  \right) \left(  \sum_{j = 1}^n g_j \partial_j (h)  \right)
}
{ =} { \sum_{i = 1}^n f_i \partial_i \left(  \sum_{j = 1}^n g_j \partial_j (h)  \right)
}
{ =} { \sum_{i = 1}^n f_i \left(  \sum_{j = 1}^n \left(  \partial_i (g_j) \partial_j (h) +  g_j \partial_i \partial_j (h)  \right)  \right)
}
{ =} { \sum_{j = 1}^n \left(  \sum_{i = 1}^n f_i  \partial_i g_j  \right) \partial_j h +  \sum_{j = 1}^n \sum_{i = 1}^n    f_i g_j \partial_i \partial_j h
}
}
{}
{}{.}

}

Dari deskripsi eksplisit ini tampak bahwa kedua komposisi pada umumnya dapat berbeda; yakni, dapat terjadi

\mavergleichskettedisp
{\vergleichskette
{ \left(  \sum_{i = 1}^n f_i \partial_i  \right) \circ  \left(  \sum_{j = 1}^n g_j \partial_j   \right)
}
{ \neq} { \left(  \sum_{j = 1}^n g_j \partial_j   \right) \circ  \left(  \sum_{i = 1}^n f_i \partial_i  \right)
}
{ } {
}
{ } {
}
{ } {
}
}
{}{}{,}
karena meskipun suku kedua selalu sama, suku pertamanya pada umumnya dapat berbeda. Hal ini juga mempunyai konsekuensi berikut.


\inputfaktbeweis
{Offene Menge/R^n/Vektorfelder/Lie-Klammer/Vektorfeld/Fakt}
{Lema}
{}
{

\faktsituation {Misalkan

\mavergleichskette
{\vergleichskette
{ U
}
{ \subseteq }{ \R^n
}
{  }{
}
{    }{
}
{   }{
}
}
{}{}{}
\definitionsverweis {terbuka}{}{}
dan misalkan $f_1 , \ldots , f_n, g_1 , \ldots , g_n$ merupakan fungsi-fungsi pada $U$ yang dua kali
\definitionsverweis {terdiferensialkan secara kontinu}{}{}
dengan operator-operator diferensial terkait

\mavergleichskette
{\vergleichskette
{ D
}
{ = }{ \sum_{i = 1}^n f_i \partial_i
}
{  }{
}
{    }{
}
{   }{
}
}
{}{}{}
dan

\mavergleichskette
{\vergleichskette
{ E
}
{ = }{ \sum_{j = 1}^n g_j \partial_j
}
{  }{
}
{    }{
}
{   }{
}
}
{}{}{}}
\faktfolgerung {Maka berlaku

\mavergleichskettedisp
{\vergleichskette
{ D \circ E -E \circ D
}
{ =} { \sum_{j = 1}^n \left(  \sum_{i = 1}^n \left(  f_i  \partial_i g_j - g_i \partial_i f_j  \right)  \right)  \partial_j
}
{ } {
}
{ } {
}
{ } {
}
}
{}{}{.}
Secara khusus, ungkapan ini sendiri bersesuaian dengan suatu operator diferensial orde $1$, dan dengan demikian dengan suatu medan vektor.}
\faktzusatz {}
\faktzusatz {}

}
{

Menurut
Lema 11.5
berlaku

\mavergleichskettealignhandlinks
{\vergleichskettealignhandlinks
{D \circ E-E \circ D
}
{ =} { \left(  \sum_{i = 1}^n f_i \partial_i  \right) \circ  \left(  \sum_{j = 1}^n g_j \partial_j   \right) -  \left(  \sum_{j = 1}^n g_j \partial_j   \right) \circ  \left(  \sum_{i = 1}^n f_i \partial_i  \right)
}
{ =} { \sum_{j = 1}^n  \left(  \sum_{i = 1}^n  \left(  f_i  \partial_i g_j - g_i \partial_i f_j  \right)  \right)  \partial_j
}
{ } {
}
{ } {
}
}
{}
{}{.}

}

Pada suatu manifold, operator diferensial yang bersesuaian dengan medan vektor $V$ kita nyatakan dengan $D_V$. Untuk suatu fungsi diferensiabel
\maabb {f} { M } {\R
} {}
operator $D_V$ didefinisikan secara titik-demi-titik oleh

\mavergleichskette
{\vergleichskette
{ \left( D_V f \right)(P)
}
{ = }{ d f_P\left(V_P\right)
}
{ \in }{ \R
}
{    }{
}
{   }{
}
}
{}{}{,}
dengan $V_P=V(P)$ dan dengan identifikasi kanonik
\mathl{T_{f(P)}\R \cong \R}{.}
Dengan kata lain, untuk suatu titik

\mavergleichskette
{\vergleichskette
{ P
}
{ \in }{ M
}
{  }{
}
{    }{
}
{   }{
}
}
{}{}{}
pemetaan singgung penuh memenuhi

\mavergleichskettedisp
{\vergleichskette
{ \left( T(f) \circ V \right)(P)
}
{ =} { \left( f(P), \left(D_V f\right)(P) \right)
}
{ } {
}
{ } {
}
{ } {
}
}
{}{}{.}
Pernyataan yang bersesuaian berlaku pula untuk fungsi diferensiabel yang didefinisikan pada suatu himpunan bagian terbuka dari $M$.

\inputdefinition
{}
{

Misalkan
\mathkor {} {V} {dan} {W} {}
merupakan
\definitionsverweis {medan-medan vektor}{}{}
yang terdiferensialkan secara kontinu pada suatu
$C^2$-\definitionsverweis {manifold diferensiabel}{}{}
dengan operator-operator diferensial terkait $D_V$ dan $D_W$. Medan vektor yang bersesuaian dengan komposisi
\mathl{D_V \circ D_W -D_W \circ D_V}{} disebut
\definitionswort {braket Lie}{}
dari kedua medan vektor tersebut. Braket ini dinotasikan dengan $[V,W]$.

}

Urutan operator yang diberi tanda minus dalam definisi merupakan suatu konvensi.


\inputfaktbeweis
{Mannigfaltigkeit/Vektorfelder/Lie-Klammer/Eigenschaften/Fakt}
{Lema}
{}
{

\faktsituation {Misalkan $M$ suatu
$C^2$-\definitionsverweis {manifold diferensiabel}{}{.}}
\faktuebergang {Maka
\definitionsverweis {braket Lie}{}{}
dari
\definitionsverweis {medan-medan vektor}{}{}
memenuhi sifat-sifat berikut.}
\faktfolgerung {\aufzaehlungdrei{Berlaku

\mavergleichskette
{\vergleichskette
{ [V,W]
}
{ = }{ -[W,V]
}
{  }{
}
{    }{
}
{   }{
}
}
{}{}{.}
}{Berlaku

\mavergleichskette
{\vergleichskette
{ [V_1+V_2,W]
}
{ = }{ [V_1,W] + [V_2,W]
}
{  }{
}
{    }{
}
{   }{
}
}
{}{}{.}
}{Untuk suatu fungsi diferensiabel $f$ pada $M$ berlaku

\mavergleichskettedisp
{\vergleichskette
{ [fV,W]
}
{ =} { f [V,W] - \left(  D_Wf  \right)  V
}
{ } {
}
{ } {
}
{ } {
}
}
{}{}{.}
}}
\faktzusatz {}
\faktzusatz {}

}
{

(1) mengikuti dari definisi, dan (2) mengikuti dari hubungan

\mavergleichskettedisp
{\vergleichskette
{ D_{V_1 +V_2 }
}
{ =} { D_{V_1 } + D_{V_2 }
}
{ } {
}
{ } {
}
{ } {
}
}
{}{}{}
antara operator diferensial dan medan vektor. (3). Karena pernyataan tersebut bersifat lokal, kita dapat menetapkan

\mavergleichskettedisp
{\vergleichskette
{ V
}
{ =} { \sum_{i = 1}^n f_i \partial_i
}
{ } {
}
{ } {
}
{ } {
}
}
{}{}{}
dan

\mavergleichskette
{\vergleichskette
{ W
}
{ = }{ \sum_{j = 1}^n g_j \partial_j
}
{  }{
}
{    }{
}
{   }{
}
}
{}{}{.}
Maka

\mavergleichskettedisp
{\vergleichskette
{ fV
}
{ =} { f \sum_{i = 1}^n f_i \partial_i
}
{ =} { \sum_{i = 1}^n ff_i \partial_i
}
{ } {
}
{ } {
}
}
{}{}{}
dan menurut
Lema 11.6
berlaku

\mavergleichskettealignhandlinks
{\vergleichskettealignhandlinks
{ \, [fV,W]
}
{ =} { \sum_{j = 1}^n \left(  \sum_{i = 1}^n \left(  f f_i \partial_i g_j - g_i \partial_i (f f_j) \right)  \right) \partial_j
}
{ =} { \sum_{j = 1}^n \left(  \sum_{i = 1}^n \left(  f f_i \partial_i g_j - g_i \left(  f \partial_i ( f_j) +f_j \partial_i (f)  \right)  \right)  \right)  \partial_j
}
{ =} { f \sum_{j = 1}^n \left(  \sum_{i = 1}^n \left(  f_i \partial_i g_j - g_i \partial_i ( f_j)  \right)  \right) \partial_j - \sum_{j = 1}^n \left(  \sum_{i = 1}^n g_i f_j \partial_i (f)  \right)  \partial_j
}
{ =} { f[V,W] - \sum_{j = 1}^n \left(  \sum_{i = 1}^n g_i \partial_i (f)  \right) f_j \partial_j
}
}
{
\vergleichskettefortsetzungalign
{ =} { f[V,W] - \left(  \sum_{i = 1}^n g_i \partial_i (f)  \right) \left(  \sum_{j = 1}^n f_j \partial_j  \right)
}
{ =} { f[V,W] - D_W(f) V
}
{ } {}
{ } {}
}
{}{.}

}
